\paragraph{R3: audited statement and notation.}
We expand Appendix A of arXiv:2407.01392v4 \cite{chen2025diffusion}. This is a corrected derivation of the intended ELBO, not a verification of every displayed formula as printed. In particular, Claim 3 prints $\mathbb E_q\log(p/q)=D_{\rm KL}(q\Vert p)$; the correct sign is negative. Theorem A.1 also reverses the arguments of the model transition in its KL. The derivation below consistently uses the reverse transition $p_\theta(x^{j-1}\mid x^j,z)$. Equation (A.13) is an indicator-expectation identity, not the Gaussian KL identity.

Let $x^0_{1:T}$ be a fixed clean sequence, with token dimension $d$, and let $k_t\in\{0,\ldots,K\}$ be independent uniform noise levels. For each token, $q(x^{1:K}\mid x^0)=\prod_{j=1}^Kq(x^j\mid x^{j-1})$ is the forward Markov chain. Write $y_t=x_t^{k_t}$ and let $z_t=z_t(k_{1:t})$ summarize the noisy history; indexing the latent by the noise history avoids assigning conflicting values to one latent variable. The model satisfies $p_\theta(y_t,z_t\mid y_{<t},z_{<t})=p_\theta(y_t,z_t\mid z_{t-1})$. Its token-level reverse chain is
\[
 p_\theta(x^{k:K}\mid z)=r(x^K)\prod_{j=k+1}^Kp_\theta(x^{j-1}\mid x^j,z),
\]
where the terminal prior $r$ is fixed and independent of $z$ and $\theta$. As in the paper, assume these conditionals arise from a consistent joint probability measure on the noise-level copies and history-indexed latents. For a fixed noise vector, $p_\theta(y_{1:T})$ denotes its selected marginal, not a mixture over noise vectors. $\mathbb E_{\rm forward}$ includes the stated draw of noise levels as well as the forward-chain variables; other expectations include the relevant chain variables even when only a marginal is displayed. Assume the densities and logarithms are integrable; otherwise the lower bound may be $-\infty$.

\paragraph{Paper Claim 1: expanding the latents.}
Define $Q_z=\prod_t p_\theta(z_t\mid z_{t-1},y_t)$. Factoring the joint distribution gives
\[
 p_\theta(y_{1:T})=\mathbb E_{Q_z}\prod_{t=1}^T p_\theta(y_t\mid z_{t-1}),\qquad
 \log p_\theta(y_{1:T})\geq\mathbb E_{Q_z}\sum_t\log p_\theta(y_t\mid z_{t-1}).
\]
\emph{Expanded step:} $Q_z$ is a normalized sequential distribution, not necessarily the full posterior conditioned on all observations. The inequality is Jensen's inequality for the concave logarithm. Deterministic latent updates make the product constant under $Q_z$, giving equality. Averaging over the forward process reduces the problem to bounding each token likelihood conditional on its sampled past.

\paragraph{Paper Claim 2: importance sampling and Jensen.}
Fix $t,k,z=z_{t-1}$, put $y=x_t^k$, $u=x_t^{k+1:K}$ and $Q(u\mid y)=q(u\mid y)$. Then
\begin{align*}
 p_\theta(y\mid z)&=\int p_\theta(y,u\mid z)\,du
 =\mathbb E_{Q(u\mid y)}\frac{p_\theta(y,u\mid z)}{Q(u\mid y)},\\
 \log p_\theta(y\mid z)&\geq
 \mathbb E_Q\left[\log p_\theta(y,u\mid z)-\log Q(u\mid y)\right]
 =:\mathcal L_k(y,z).
\end{align*}
\emph{Expanded step:} the importance-sampling equality requires the model measure on $u$ to be absolutely continuous with respect to $Q$: no model mass may lie where $Q=0$. Nondegenerate Gaussian chains satisfy this. A finite ELBO also requires $Q$ not to put positive mass where the model density is zero. Concavity gives $\log\mathbb E_Q W\geq\mathbb E_Q\log W$. More explicitly,
\[
 \log p_\theta(y\mid z)-\mathcal L_k(y,z)
 =D_{\rm KL}\bigl(Q(u\mid y)\Vert p_\theta(u\mid y,z)\bigr)\geq0.
\]
Thus equality requires an exact variational posterior. When $k=K$, $u$ is empty and this step is an equality. These are the paper's two Jensen inequalities, including the latent step above.

\paragraph{Paper Claim 3: time reversal and telescoping.}
For $k\geq1$, Bayes' rule and the forward Markov property give
\[
 q(x^j\mid x^{j-1})=
 q(x^{j-1}\mid x^j,x^0)\frac{q(x^j\mid x^0)}{q(x^{j-1}\mid x^0)}.
\]
\emph{Expanded step:} this follows by expressing the same conditional joint density of $(x^{j-1},x^j)$ given $x^0$ in both orders. Substituting it into the log likelihood ratio produces the marginal-density sum
\begin{align*}
 &\bigl[\log q(x^k\mid x^0)-\log q(x^{k+1}\mid x^0)\bigr]\\
 &+\bigl[\log q(x^{k+1}\mid x^0)-\log q(x^{k+2}\mid x^0)\bigr]+\cdots\\
 &+\bigl[\log q(x^{K-1}\mid x^0)-\log q(x^K\mid x^0)\bigr]\\
 &=\log q(x^k\mid x^0)-\log q(x^K\mid x^0).
\end{align*}
Every intermediate marginal cancels. Consequently, after taking the forward expectation,
\[
 \mathbb E_q\mathcal L_k=C_k(x^0)
 -\sum_{j=k+1}^K\mathbb E_{q(x^j\mid x^0)}
 D_{\rm KL}\bigl(q(x^{j-1}\mid x^j,x^0)\Vert
 p_\theta(x^{j-1}\mid x^j,z)\bigr),\quad k\geq1,
\]
where $C_k=\mathbb E_q[\log r(x^K)-\log q(x^K\mid x^0)+\log q(x^k\mid x^0)]$ is parameter independent. This also handles $k=K$ with an empty sum. For $k=0$, the clean endpoint is fixed, so one must keep the reconstruction term instead of using a density for a Dirac mass:
\[
 \mathbb E_q\mathcal L_0=C_0(x^0)+\mathbb E_q\log p_\theta(x^0\mid x^1,z)
 -\sum_{j=2}^K\mathbb E_{q(x^j\mid x^0)}D_{\rm KL}(q(x^{j-1}\mid x^j,x^0)\Vert p_\theta(x^{j-1}\mid x^j,z)),
\]
with $C_0=\mathbb E_q[\log r(x^K)-\log q(x^K\mid x^0)]$. These signs follow directly from $\int q\log(p/q)=-D_{\rm KL}(q\Vert p)$.

\paragraph{Equation (A.13), expanded.}
There is one choice $k=0$ out of $K+1$ possibilities. For $0\leq j<K$, the event $k\leq j$ contains $j+1$ possibilities, so
\[
 \mathbb E\mathbf1\{k=0\}=\frac1{K+1},\qquad
 \mathbb E\mathbf1\{j\geq k\}=\sum_{k=0}^j\frac1{K+1}=\frac{j+1}{K+1}.
\]
A reverse transition from $j$ to $j-1$ appears whenever $k<j$, hence has weight $j/(K+1)$. Combining the three claims yields the corrected intended statement
\begin{align*}
 \mathbb E_{\rm forward}\log p_\theta(y_{1:T})\geq C+
 \sum_t\mathbb E_{<t}\biggl[&\frac{1}{K+1}\mathbb E_q\log p_\theta(x_t^0\mid x_t^1,z_{t-1})\\
 &-\sum_{j=2}^K\frac{j}{K+1}\mathbb E_{q(x_t^j\mid x_t^0)}D_{t,j}\biggr],
\end{align*}
where $D_{t,j}$ is the reverse-transition KL above and $\mathbb E_{<t}$ averages noisy history and latent updates. In contrast to the printed positive-KL expression, improving a denoiser increases this lower bound.

\paragraph{Expanded Gaussian KL reduction.}
Let $q(x^j\mid x^{j-1})=\mathcal N(\sqrt{\alpha_j}x^{j-1},\beta_j I)$, $\alpha_j=1-\beta_j$, and $\bar\alpha_j=\prod_{i=1}^j\alpha_i$. Completing the square in
$q(x^j\mid x^{j-1})q(x^{j-1}\mid x^0)$ gives posterior variance and mean
\[
 \widetilde\beta_j=\frac{\beta_j(1-\bar\alpha_{j-1})}{1-\bar\alpha_j},\qquad
 \widetilde\mu_j=A_jx^j+B_jx^0,
\]
\[
 A_j=\frac{\sqrt{\alpha_j}(1-\bar\alpha_{j-1})}{1-\bar\alpha_j},\qquad
 B_j=\frac{\beta_j\sqrt{\bar\alpha_{j-1}}}{1-\bar\alpha_j}.
\]
Indeed, the posterior precision is $\alpha_j/\beta_j+1/(1-\bar\alpha_{j-1})=1/\widetilde\beta_j$; multiplying its inverse by the linear coefficient gives the displayed mean. For $j\geq2$, take model covariance $\sigma_j^2I>0$, fixed independently of $\theta$, and model mean $A_jx^j+B_j\widehat x^0_\theta$. The Gaussian KL is
\[
 D_{t,j}=\frac12\left[d\left(\frac{\widetilde\beta_j}{\sigma_j^2}-1+\log\frac{\sigma_j^2}{\widetilde\beta_j}\right)
 +\frac{B_j^2}{\sigma_j^2}\|\widehat x^0_\theta-x^0\|^2\right].
\]
The first term is parameter independent. Thus $c_j=B_j^2/(2\sigma_j^2)$ multiplies the squared error, and $j/(K+1)$ supplies the additional partial-noise weight. At $j=1$, use a separate finite-variance Gaussian decoder; its negative log likelihood is a constant plus $\|\widehat x^0-x^0\|^2/(2\sigma_1^2)$. Setting its variance to the degenerate forward posterior variance would be invalid.

Writing $x^j=\sqrt{\bar\alpha_j}x^0+\sqrt{1-\bar\alpha_j}\epsilon$, consistent parameterizations give
\[
 \|\widehat x^0-x^0\|^2=\frac{1-\bar\alpha_j}{\bar\alpha_j}\|\widehat\epsilon-\epsilon\|^2
 =(1-\bar\alpha_j)\|\widehat v-v\|^2,
\]
where $v=\sqrt{\bar\alpha_j}\epsilon-\sqrt{1-\bar\alpha_j}x^0$ and $x^0=\sqrt{\bar\alpha_j}x^j-\sqrt{1-\bar\alpha_j}v$. These identities supply the corresponding loss weights without changing the ELBO sign.

\paragraph{Expanded subsequence marginalization.}
Let $S\subseteq\{1,\ldots,T\}$ be the observed token indices. Set $k_t=0$ for $t\in S$ and $k_t=K$ otherwise. Under exact terminal masking, the omitted forward tokens $u=(x_t^K)_{t\notin S}$ have a data-independent density $r_S(u)=\prod_{t\notin S}r(x_t^K)$. For the consistent joint measure assumed above,
\begin{align*}
\log p_\theta(x_S^0)
 &=\log\mathbb E_{r_S}\frac{p_\theta(x_S^0,u)}{r_S(u)}\\
 &\geq\mathbb E_{r_S}\log p_\theta(x_S^0,u)-\mathbb E_{r_S}\log r_S(u).
\end{align*}
The importance-sampling support condition is the same as before. This additional Jensen step has gap
$D_{\rm KL}(r_S(u)\Vert p_\theta(u\mid x_S^0))$; the noise-entropy term is independent of $\theta$. Applying the fixed-noise-vector bound to its right-hand side gives an ELBO for the subsequence marginal. Observed tokens retain every denoising level, while fully masked tokens contribute no reverse-transition loss. These token/history combinations are covered by independent-noise training, up to the sufficient-capacity and endpoint qualifications below. For $S=\{1,\ldots,T\}$ the omitted variables are empty, recovering the full clean-sequence case. This argument assumes consistent marginals; an arbitrary collection of separately specified conditional networks does not establish that assumption on its own.

\paragraph{Subsequences and limits of the guarantee.}
For an arbitrary distribution of noise histories, the coefficient on reverse step $j$ is $\Pr(k_t<j\mid k_{<t})$, and the decoder coefficient is $\Pr(k_t=0\mid k_{<t})$. They are cumulative probabilities, not generally $j\Pr(k_t=j\mid k_{<t})$ as printed in Theorem A.3. Clean tokens and completely masked tokens represent subsequence observations when terminal noise is independent of clean data; finite Gaussian schedules with nonzero terminal signal only approximate complete masking. Independent-noise training covers these histories. Dropping the weights preserves simultaneous optima only under the paper's idealized sufficient-capacity assumptions; it need not do so for a finite shared network. The bound is on expected log likelihood of partially noised sequences, not a proof of rollout stability or of tightness for a trained checkpoint. The released base configuration uses \texttt{noise\_level: random\_all}; its implementation draws indices from $0$ through $999$ for $1000$ diffusion steps. Code index zero is the first positive-noise step, rather than the theorem's exactly clean endpoint $k=0$. Video additionally uses fused-SNR weighting and a finite attention window. Our inference-only audit therefore establishes a connection to the training principle, not empirical verification of the theorem's idealized assumptions.
